\chapter{Twistory: spinory, incydencja i fibracja nad sferą}
\label{ch:lin32-twistory}
\index{twistor}
Rozdział~\ref{ch:lin31-spinory} wprowadził spinory i ich reprezentacje.
Teraz pokażemy, jak algebra liniowa spinorów prowadzi do geometrii
twistorowej. W modelu zespolonym punkt przestrzeni czterowymiarowej
odpowiada prostej rzutowej w $\mathbb{CP}^3$. W modelu euklidesowym
nad każdym punktem $S^4$ znajduje się taka prosta, będąca włóknem
gładkiej fibracji. Konstrukcje te mają tę samą przestrzeń rzutową,
lecz korzystają z różnych dodatkowych struktur. Rozróżnienie
obu modeli jest istotne.

\section{Macierze $2\times2$ jako wektory czterowymiarowe}
\label{sec:lin32-macierze}
Ustalmy dwie dwuwymiarowe zespolone przestrzenie $S_+$ i $S_-$
z wybranymi bazami. Ich przestrzeń homomorfizmów
$V_{\C}:=\operatorname{Hom}_{\C}(S_-,S_+)$
utożsamiamy z $M_2(\C)$. Odpowiada to dwóm typom
spinorów chiralnych z Rozdziału~\ref{ch:lin31-spinory},
lecz samo utożsamienie wymaga wyboru baz.

\begin{definicja}[Forma kwadratowa wyznacznika]
\label{def:lin32-forma}
Dla $X=\begin{psmallmatrix}a&b\\c&d\end{psmallmatrix}$
określamy $q(X)=\det X=ad-bc$ i jego polaryzację
$B(X,Y)=q(X+Y)-q(X)-q(Y)$.
Niezerową macierz $X$ nazywamy \emph{izotropową},
gdy $q(X)=0$. Punkty $X,Y$ nazywamy
\emph{rozdzielonymi izotropowo}, gdy $q(X-Y)=0$.
\end{definicja}

\begin{stwierdzenie}[Niezdegenerowanie i działanie grup]
\label{prop:lin32-dzialanie}
Forma $B$ jest niezdegenerowana.
Grupa $\mathrm{SL}_2(\C)\times\mathrm{SL}_2(\C)$ działa na
$V_{\C}$ przez $(A,D)\cdot X=AXD^{-1}$, zachowując $q$.
Jądro tego działania to
$\{(I_2,I_2),(-I_2,-I_2)\}$.
\end{stwierdzenie}
\begin{proof}
Dla $Y=\begin{psmallmatrix}a'&b'\\c'&d'\end{psmallmatrix}$
rozwinięcie wyznacznika daje
$B(X,Y)=ad'+a'd-bc'-b'c$.
Jeśli $B(X,Y)=0$ dla wszystkich $Y$, podstawienie
czterech macierzy jednostkowych $E_{11},E_{22},E_{12},E_{21}$
wymusza kolejno $d=a=c=b=0$.
Z multiplikatywności wyznacznika
$q(AXD^{-1})=(\det A)q(X)(\det D)^{-1}=q(X)$.
Jeśli $(A,D)$ działa trywialnie, to z $X=I_2$
wynika $A=D$. Wtedy $AX=XA$ dla wszystkich $X$.
Komutowanie z $E_{11}$ wymusza diagonalność $A$,
a z $E_{12}$ równość wyrazów diagonalnych.
Stąd $A=\lambda I_2$ i $\lambda^2=\det A=1$.
Obie pary znaków rzeczywiście działają trywialnie.
\end{proof}

\begin{przyklad}[Postać rzeczywista lorentzowska]
\label{prz:lin32-hermit}
Dla rzeczywistych $t,x,y,z$ bierzemy macierz hermitowską
\[
 H(t,x,y,z)=
 \begin{pmatrix}t+z&x-iy\\x+iy&t-z\end{pmatrix}.
\]
Mnożenie liczb zespolonych daje
$\det H=(t+z)(t-z)-(x-iy)(x+iy)=t^2-x^2-y^2-z^2$.
Jest to rzeczywista forma o sygnaturze $(1,3)$.
Nie jest ona euklidesową metryką na $\R^4$:
tę odmienną sytuację rozpatrzymy przy fibracji nad $S^4$.
\end{przyklad}

\section{Kwadryka Segrego i płaszczyzny izotropowe}
\label{sec:lin32-segre}
Każda niezerowa macierz $X$ o $\det X=0$ ma rząd
jeden, więc można ją zapisać jako
$X=u\alpha^{\mathsf T}$ dla niezerowych
$u,\alpha\in\C^2$. Dwa takie przedstawienia
różnią się zastąpieniem $(u,\alpha)$ przez
$(\lambda u,\lambda^{-1}\alpha)$.

\begin{twierdzenie}[Kwadryka Segrego]
\label{thm:lin32-segre}
Odwzorowanie
\[
 \mathrm{Seg}:\mathbb{CP}^1\times\mathbb{CP}^1
 \longrightarrow\mathbb{CP}^3,\qquad
 ([u_0:u_1],[a_0:a_1])
 \longmapsto[u_0a_0:u_0a_1:u_1a_0:u_1a_1]
\]
jest osadzeniem holomorficznym. Jego obrazem jest
kwadryka $Q=\{[z_{11}:z_{12}:z_{21}:z_{22}]:
z_{11}z_{22}-z_{12}z_{21}=0\}$, składająca się
dokładnie z rzutowych klas niezerowych macierzy
izotropowych w $V_{\C}$.
\end{twierdzenie}
\begin{proof}
Skalowanie któregokolwiek czynnika mnoży wszystkie
współrzędne obrazu przez wspólną niezerową liczbę,
więc wzór jest dobrze określony. Jego macierz
ma rząd jeden, a zatem wyznacznik zero.
Na odwrót, każda niezerowa macierz o wyznaczniku
zero ma jednoznacznie określoną prostą obrazu
i prostą przestrzeni wierszy. Te dwie proste
odtwarzają $[u]$ i $[\alpha]$, więc odwzorowanie
jest bijekcją na $Q$.
Na mapach $u_0a_0\ne0$ piszemy $u=(1,s)$,
$\alpha=(1,t)$; wtedy obrazem jest $[1:t:s:st]$.
Jego odwrotność w tej mapie wyznaczają
$s=z_{21}/z_{11}$ i $t=z_{12}/z_{11}$.
Pozostałe trzy pary map otrzymujemy przez
zmianę wyróżnionych niezerowych współrzędnych.
W każdej z nich odwrotność jest holomorficzna,
co dowodzi osadzenia.
\end{proof}

\begin{twierdzenie}[Dwie rodziny maksymalnych płaszczyzn izotropowych]
\label{thm:lin32-rodziny}
Każda dwuwymiarowa podprzestrzeń $L\subset V_{\C}$
spełniająca $q|_L=0$ ma postać
\[
 \{u\alpha^{\mathsf T}:\alpha\in\C^2\}
 \quad\text{lub}\quad
 \{u\alpha^{\mathsf T}:u\in\C^2\},
\]
dla ustalonego niezerowego $u$ albo ustalonego
niezerowego $\alpha$. Są to wszystkie maksymalne
podprzestrzenie izotropowe, parametryzowane
przez dwie kopie $\mathbb{CP}^1$.
\end{twierdzenie}
\begin{proof}
Wybierzmy niezależne
$X=u\alpha^{\mathsf T}$, $Y=v\beta^{\mathsf T}$
rozpinające $L$. Dla każdego $\lambda\in\C$
mamy
\[
 0=\det(X+\lambda Y)
 =\lambda\,\det[u\ v]\,
 \det\begin{pmatrix}\alpha^{\mathsf T}\\
                     \beta^{\mathsf T}\end{pmatrix}.
\]
Zatem albo $u,v$ są proporcjonalne, albo
$\alpha,\beta$ są proporcjonalne.
W pierwszym przypadku wszystkie macierze z $L$
mają tę samą prostą obrazu, a w drugim
tę samą przestrzeń wierszy. Otrzymujemy dwie postacie
z tezy; obie rzeczywiście są podprzestrzeniami
wymiaru dwa i zawierają tylko macierze rzędu
co najwyżej jeden. Jeśli $W$ jest dowolną
izotropową podprzestrzenią, polaryzacja
$q(w+w')=0$ pokazuje, że $B|_{W\times W}=0$,
więc $W\subset W^{\perp_B}$.
Niezdegenerowanie $B$ daje
$\dim W+\dim W^{\perp_B}=4$,
czyli $\dim W\le2$. Udowodniliśmy maksymalność.
\end{proof}

Rzutowanie takich płaszczyzn daje dwie rodziny
prostych leżących na kwadryce $Q$.
W teorii twistorów prowadzi to do dwóch rodzajów
całkowicie izotropowych płaszczyzn, zwanych
często płaszczyznami $\alpha$ i $\beta$.
Z konwencją incydencji poniżej związana będzie
rodzina o ustalonej przestrzeni wierszy.

\section{Twistory i incydencja}
\label{sec:lin32-incydencja}
\begin{definicja}[Przestrzeń twistorów i prosta twistorowa]
\label{def:lin32-twistor}
Niech $\mathbb T=S_+\oplus S_-\cong\C^4$.
Niezerowy \emph{twistor} zapisujemy
$Z=(\omega,\pi)$, $\omega\in S_+$, $\pi\in S_-$.
\emph{Rzutowa przestrzeń twistorów} to
$\mathbb{PT}=\mathbb P(\mathbb T)\cong\mathbb{CP}^3$.
Dla punktu $X\in V_{\C}=M_2(\C)$
\emph{relacja incydencji} ma postać
$\omega=X\pi$ przy $\pi\ne0$.
Odpowiadająca $X$ \emph{prosta twistorowa} to
\[
 \ell_X=\{[X\pi:\pi]:\pi\in\C^2\setminus\{0\}\}
                   \subset\mathbb{PT}.
\]
Konwencja fizyczna $\omega=iX\pi$ powstaje
po zastąpieniu $X$ przez $iX$.
\end{definicja}

Przestrzeń $\{(X\pi,\pi):\pi\in\C^2\}$ jest
dwuwymiarową zespoloną podprzestrzenią $\C^4$,
będącą wykresem odwzorowania liniowego $X$.
Jej rzutowanie jest prostą $\mathbb{CP}^1$.
Dlatego punkt $X$ odpowiada \emph{całej prostej},
a nie jednemu twistorowi.

\begin{twierdzenie}[Przecięcie prostych i izotropowość]
\label{thm:lin32-przeciecie}
Dla dowolnych $X,Y\in M_2(\C)$ zachodzi
\[
 \ell_X\cap\ell_Y\ne\varnothing
 \quad\Longleftrightarrow\quad\det(X-Y)=0.
\]
Dla $X\ne Y$ z zerowym wyznacznikiem różnicy
proste przecinają się w dokładnie jednym punkcie;
dla $X=Y$ pokrywają się.
\end{twierdzenie}
\begin{proof}
Punkt wspólny ma przedstawienia
$[X\pi:\pi]=[Y\eta:\eta]$ z $\pi,\eta\ne0$.
Istnieje więc $\lambda\in\C^\ast$ takie, że
$X\pi=\lambda Y\eta$ i $\pi=\lambda\eta$,
czyli $(X-Y)\pi=0$.
Na odwrót, każdy niezerowy $\pi\in\ker(X-Y)$
daje punkt obu prostych.
Jądro jest niezerowe dokładnie przy
$\det(X-Y)=0$. Gdy $X\ne Y$, macierz różnicy
jest niezerowa o rzędzie jeden, a jej jądro
ma wymiar jeden, więc rzutowo otrzymujemy
dokładnie jeden punkt. Przypadek $X=Y$
jest bezpośredni.
\end{proof}

\begin{przyklad}[Przecięcie i rozłączność]
\label{prz:lin32-przeciecie}
Dla $X=0$ oraz $Y=\operatorname{diag}(1,0)$
różnica ma wyznacznik zero i
$\ker(X-Y)=\C(0,1)^{\mathsf T}$.
Wspólnym twistorem jest $[0:0:0:1]$.
Natomiast dla $X=0$ oraz $Y=I_2$
wyznacznik różnicy wynosi $1$,
więc proste nie mają punktu wspólnego.
\end{przyklad}

\begin{stwierdzenie}[Płaszczyzna incydencji jednego twistoru]
\label{prop:lin32-alfa}
Dla $Z=[\omega:\pi]\in\mathbb{PT}$,
$\pi\ne0$, zbiór
$A_Z=\{X\in M_2(\C):X\pi=\omega\}$
jest afiniczną płaszczyzną zespoloną
wymiaru dwa. Jej przestrzeń kierunkowa
$N_\pi=\{T:T\pi=0\}$ jest całkowicie izotropowa
i należy do rodziny o stałej przestrzeni
wierszy z Twierdzenia~\ref{thm:lin32-rodziny}.
\end{stwierdzenie}
\begin{proof}
Odwzorowanie $L_\pi:M_2(\C)\to\C^2$,
$L_\pi(X)=X\pi$, jest surjektywne:
wybierając $\eta$ tak, aby $(\pi,\eta)$
stanowiło bazę, możemy zadawać dowolnie
$X\pi=\omega$ i $X\eta=0$.
Z twierdzenia o rzędzie i jądrze
$\dim_\C N_\pi=4-2=2$.
Jego każde niezerowe $T$ ma niezerowe
jądro, więc $\det T=0$. Ta sama własność
dotyczy kombinacji liniowych elementów $N_\pi$.
Jeśli niezerowy funkcjonał $\alpha$ ma
$\ker\alpha=\C\pi$, każdy wiersz $T$
jest proporcjonalny do $\alpha$.
Stąd $N_\pi=\{u\alpha:u\in\C^2\}$.
\end{proof}

\begin{przyklad}[Jawna płaszczyzna $\alpha$]
\label{prz:lin32-alfa}
Dla $\pi=(1,0)^{\mathsf T}$,
$\omega=(a,b)^{\mathsf T}$ warunek
$X\pi=\omega$ ustala tylko pierwszą kolumnę:
\[
 A_Z=\left\{\begin{pmatrix}a&s\\b&t\end{pmatrix}
                 :s,t\in\C\right\}.
\]
Każda różnica dwóch takich macierzy
ma pierwszą kolumnę zerową i wyznacznik zero.
\end{przyklad}

\begin{figure}[!htbp]
\centering
\begin{tikzpicture}[>=Stealth,font=\small]
 \node[draw,rounded corners,fill=manifoldblue!8,
       minimum width=3.2cm,minimum height=1.0cm]
       (F) at (0,2.1)
       {$\mathcal F_{\mathrm{aff}}=M_2(\C)\times\mathbb{CP}^1$};
 \node[draw,rounded corners,fill=manifoldgreen!9,
       minimum width=2.5cm,minimum height=.9cm]
       (M) at (-3.1,0) {$M_2(\C)$};
 \node[draw,rounded corners,fill=manifoldgold!12,
       minimum width=2.8cm,minimum height=.9cm]
       (P) at (3.1,0) {$\mathbb{CP}^3$};
 \draw[->,thick,manifoldblue] (F.south west)--(M.north)
       node[midway,left] {$(X,[\pi])\mapsto X$};
 \draw[->,thick,manifoldblue] (F.south east)--(P.north)
       node[midway,right] {$[X\pi:\pi]$};
 \node[align=center] at (0,-1.05)
       {Nad $X$: prosta $\ell_X$.\\
        Nad $[\omega:\pi]$: płaszczyzna $A_Z$.};
\end{tikzpicture}
\caption{Dwa rzuty korespondencji incydencji.
Punkt przestrzeni macierzy odpowiada prostej twistorowej,
a twistor z $\pi\ne0$ płaszczyźnie afinicznej.}
\label{fig:lin32-incydencja}
\end{figure}

\section{Grassmannian i kwadryka Kleina}
\label{sec:lin32-klein}
Jak w Rozdziale~\ref{ch:lin29-grassmanniany},
$\operatorname{Gr}_2(\C^4)$ parametryzuje
dwuwymiarowe podprzestrzenie $\C^4$.
Punktowi $X$ przypisujemy wykres
$W_X=\{(X\pi,\pi):\pi\in\C^2\}$.
Rzutowanie $W_X\to S_-$ jest izomorfizmem.
Odwrotnie, każda płaszczyzna o takim rzutowaniu
jest wykresem jednoznacznego $X$,
więc $M_2(\C)$ jest otwartą mapą Grassmannianu.

\begin{twierdzenie}[Współrzędne Plückera wykresu]
\label{thm:lin32-plucker}
Przy osadzeniu Plückera
$\operatorname{Gr}_2(\C^4)\hookrightarrow\mathbb{CP}^5$
obrazem jest kwadryka Kleina
\begin{equation}\label{eq:lin32-klein}
 p_{12}p_{34}-p_{13}p_{24}+p_{14}p_{23}=0.
\end{equation}
Dla $X=\begin{psmallmatrix}a&b\\c&d\end{psmallmatrix}$
punkt $W_X$ ma współrzędne
\[
 [p_{12}:p_{13}:p_{14}:p_{23}:p_{24}:p_{34}]
       =[ad-bc:-b:a:-d:c:1].
\]
Część $p_{34}\ne0$ tej kwadryki jest
izomorficzna z $M_2(\C)$.
\end{twierdzenie}
\begin{proof}
Bazą $W_X$ są kolumny
$v_1=(a,c,1,0)^{\mathsf T}$,
$v_2=(b,d,0,1)^{\mathsf T}$.
Wyznaczniki dwuwierszowych minorów
dają kolejno $ad-bc,-b,a,-d,c,1$.
Podstawienie w~\eqref{eq:lin32-klein}
daje $(ad-bc)+bc-ad=0$.
Dla ogólnej płaszczyzny równanie
Plückera zostało dowiedzione w
Twierdzeniu~\ref{thm:lin29-grass-plucker}.
Jeśli $p_{34}\ne0$, dolne dwie
współrzędne pewnej bazy płaszczyzny
tworzą macierz odwracalną.
Zmieniając bazę, sprowadzamy jej dolny
blok do $I_2$ i otrzymujemy dokładnie
płaszczyznę $W_X$ dla jednej macierzy $X$.
\end{proof}

Nie należy mylić dwóch kwadryk:
$Q\subset\mathbb{CP}^3$ z
Twierdzenia~\ref{thm:lin32-segre}
opisuje \emph{izotropowe kierunki}
w $V_{\C}$, podczas gdy czterowymiarowa
kwadryka Kleina w $\mathbb{CP}^5$,
czyli $\operatorname{Gr}_2(\C^4)$,
opisuje rzutowe domknięcie
\emph{przestrzeni punktów} $X$.

\begin{stwierdzenie}[Pełna korespondencja flag]
\label{prop:lin32-flagi}
Niech $\mathcal F$ będzie zbiorem par
$(\ell,W)$ podprzestrzeni $\C^4$
takich, że $\dim\ell=1$, $\dim W=2$,
$\ell\subset W$. Rzut na
$\operatorname{Gr}_2(\C^4)$ ma włókno
$\mathbb{CP}^1$, a rzut na
$\operatorname{Gr}_1(\C^4)=\mathbb{CP}^3$
ma włókno $\mathbb{CP}^2$.
\end{stwierdzenie}
\begin{proof}
Nad ustalonym $W$ wybieramy prostą
$\ell\subset W$, więc włóknem jest
$\mathbb P(W)\cong\mathbb{CP}^1$.
Nad ustalonym $\ell$ każda
płaszczyzna $W\supset\ell$ odpowiada
prostej $W/\ell\subset\C^4/\ell$.
Ponieważ przestrzeń ilorazowa ma
wymiar trzy, jej proste tworzą
$\mathbb{CP}^2$.
Dla $W=W_X$ i $\ell=[\omega:\pi]$
otrzymujemy ponownie incydencję
$\omega=X\pi$, o ile $\pi\ne0$.
\end{proof}

\section{Kwaternionowa struktura na $\C^4$}
\label{sec:lin32-j}
W drugim modelu $\mathbb{CP}^3$ wybieramy
\emph{rzeczywistą strukturę kwaternionową}.
Zapiszmy $\mathbb H=\{a+jb:a,b\in\C\}$,
z $j^2=-1$ i $aj=j\bar a$.
W prawostronnej przestrzeni zespolonej
$\mathbb H\cong\C^2$ prawe mnożenie przez $j$
ma współrzędne
\[
 J_0(a,b)=(-\bar b,\bar a),\qquad
 J_0^2=-\id.
\]
Ostatni wzór wynika po podstawieniu
$J_0(-\bar b,\bar a)=(-a,-b)$.
Na $\C^4\cong\mathbb H^2$ ustalamy
$J=J_0\oplus J_0$.

\begin{lemat}[Rzutowa inwolucja bez punktów stałych]
\label{lem:lin32-rzeczywista}
Dla $z\in\C^4\setminus\{0\}$ wektory
$z,Jz$ są liniowo niezależne,
a $W_z=\operatorname{span}_\C(z,Jz)$
jest $J$-niezmienniczą płaszczyzną.
Odwzorowanie $\jmath([z])=[Jz]$
jest antyholomorficzną inwolucją
$\mathbb{CP}^3$ bez punktów stałych.
\end{lemat}
\begin{proof}
Gdyby $Jz=\lambda z$, to
$-z=J^2z=\bar\lambda Jz=|\lambda|^2z$,
sprzecznie z $z\ne0$.
Niezależność jest więc dowiedziona,
a $J(W_z)=W_z$ wynika z $J(Jz)=-z$.
Dla $\mu\ne0$ mamy
$J(\mu z)=\bar\mu Jz$, więc
$\jmath$ nie zależy od przedstawiciela
projektowego i jest antyholomorficzne.
Równość $[J^2z]=[-z]=[z]$
daje inwolucję, a jej punkt stały
prowadziłby do wykluczonej równości
$Jz=\lambda z$.
\end{proof}

Każda $J$-niezmiennicza zespolona
płaszczyzna wymiaru dwa jest
jednowymiarową prawostronną
podprzestrzenią kwaternionową
w $\mathbb H^2$. Istotnie, jest
zamknięta na mnożenie przez
$\C$ i przez $j$, a więc przez
wszystkie kwaterniony; dowolny
niezerowy $z$ wraz z $Jz$
stanowi jej zespoloną bazę.
Rodzinę tych płaszczyzn oznaczamy
$\mathbb{HP}^1$.

\begin{twierdzenie}[Fibracja twistorowa]
\label{thm:lin32-fibracja}
Odwzorowanie
\begin{equation}\label{eq:lin32-fibracja}
 \tau:\mathbb{CP}^3\longrightarrow\mathbb{HP}^1,
 \qquad \tau([z])=\operatorname{span}_\C(z,Jz)
\end{equation}
jest gładką, surjektywną, lokalnie trywialną
fibracją z włóknem $\mathbb{CP}^1$.
Ponadto $\mathbb{HP}^1$ jest
dyfeomorficzna z $S^4$.
\end{twierdzenie}
\begin{proof}
Z $J(\lambda z)=\bar\lambda Jz$
wynika niezależność od przedstawiciela.
Każda kwaternionowa prosta $W$
ma bazę zespoloną $(z,Jz)$,
więc należy do obrazu $\tau$.
Przeciwobraz $W$ składa się ze wszystkich
zespolonych prostych w $W$,
czyli jest $\mathbb P(W)\cong\mathbb{CP}^1$.

Udowodnijmy gładką lokalną trywialność.
Punkty $\mathbb{HP}^1$ zapisujemy
jako $[q_0:q_1]_{\mathbb H}$, z
$(q_0,q_1)\sim(q_0h,q_1h)$
dla $h\in\mathbb H^\ast$.
W obszarze $q_0\ne0$ współrzędną jest
$q=q_1q_0^{-1}\in\mathbb H$,
a odpowiadająca prosta ma postać
$W_q=\{(v,qv):v\in\mathbb H\}$.
Wybór zespolonej prostej $\C v\subset\mathbb H$
daje trywializację
\[
 (q,[v]_\C)\longmapsto[(v,qv)]_\C
 \quad\text{z }\mathbb H\times\mathbb{CP}^1
 \text{ na }\tau^{-1}(\{q_0\ne0\}).
\]
Wzór jest bijektywny: pierwsza składowa
niezerowego wektora w $W_q$ nie może
być zerem, a jej klasa wyznacza $[v]_\C$.
Dla $q=s+jt$, $v=a+jb$ mnożenie
kwaternionów daje
$qv=(sa-\bar t b)+j(ta+\bar s b)$.
Zatem trywializacja ma współrzędne
\[
 (q,[a:b])\longmapsto
 [a:b:sa-\bar t b:ta+\bar s b].
\]
Jej odwrotność otrzymujemy przez
$q=q_1q_0^{-1}$ i $[v]_\C=[q_0]_\C$
po utożsamieniu pary zespolonej z kwaternionem.
Oba wzory są gładkie w każdej mapie
rzutowej. Nad $q_1\ne0$ postępujemy
symetrycznie.

Pozostaje zidentyfikować podstawę.
Obszar $q_0\ne0$ daje $\mathbb H\cong\R^4$,
a jego dopełnieniem jest jeden punkt
$\infty=[0:1]_{\mathbb H}$.
Dla $q\in\mathbb H$ połóżmy
\[
 \Phi(q)=
 \left(\frac{2q}{1+|q|^2},
       \frac{|q|^2-1}{1+|q|^2}\right)
       \in\mathbb H\oplus\R\cong\R^5.
\]
To punkt $S^4$, ponieważ
$4|q|^2+(|q|^2-1)^2=(1+|q|^2)^2$.
Przy $q\to\infty$ mamy
$\Phi(q)\to(0,1)$.
W drugiej mapie kwaternionowej
współrzędną na przecięciu jest $q^{-1}$;
wzór w tych współrzędnych przedłuża się
gładko przez punkt $\infty$.
Otrzymujemy stereograficzną
dyfeomorficzną identyfikację
$\mathbb{HP}^1\cong S^4$.
\end{proof}

\begin{przyklad}[Dwa jawne włókna]
\label{prz:lin32-wlokna}
Nad punktem $q=0$ mamy
$W_0=\{(v,0):v\in\mathbb H\}$,
więc
$\tau^{-1}(0)=\{[z_0:z_1:0:0]\}\cong\mathbb{CP}^1$.
Nad $\infty$ mamy
$\tau^{-1}(\infty)=\{[0:0:z_2:z_3]\}\cong\mathbb{CP}^1$.
Te proste są rozłączne. Ogólniej,
dwa różne włókna są rozłączne,
ponieważ różne jednowymiarowe
podprzestrzenie kwaternionowe
mają przecięcie $\{0\}$.
\end{przyklad}

\begin{figure}[!htbp]
\centering
\begin{tikzpicture}[>=Stealth,font=\small]
 \draw[manifoldblue!75!black,thick] (0,2.25)
       ellipse (2.65 and .85);
 \node at (-.95,2.55) {$\mathbb{CP}^3$};
 \draw[manifoldgreen!65!black,thick] (.85,2.23)
       ellipse (.78 and .34);
 \node[manifoldgreen!55!black] at (.85,2.23)
       {$\mathbb{CP}^1$};
 \draw[->,thick,manifoldgold!85!black]
       (0,1.36)--(0,.12) node[midway,right] {$\tau$};
 \draw[manifoldblue!75!black,thick] (0,-.84) circle (.78);
 \draw[manifoldblue!50] (-.78,-.84)
       arc[start angle=180,end angle=360,
       x radius=.78,y radius=.18];
 \draw[manifoldblue!50,dashed] (.78,-.84)
       arc[start angle=0,end angle=180,
       x radius=.78,y radius=.18];
 \fill[manifoldred] (.16,-.35) circle (2pt)
       node[right] {$x$};
 \node at (1.23,-1.13) {$S^4$};
\end{tikzpicture}
\caption{Schemat fibracji twistorowej
$\mathbb{CP}^1\hookrightarrow\mathbb{CP}^3
\xrightarrow{\tau}S^4$. Rysunek nie oznacza,
że całkowita przestrzeń jest globalnym iloczynem.}
\label{fig:lin32-fibracja}
\end{figure}

\begin{uwaga}[Dwa modele, dwie geometrie]
\label{rem:lin32-modele}
Zespolona relacja incydencji pozwala, by
różne proste $\ell_X,\ell_Y$ miały punkt wspólny;
dzieje się tak właśnie wtedy, gdy
$\det(X-Y)=0$. W modelu euklidesowym
różne włókna $\tau^{-1}(x)$ są rozłączne.
Nie utożsamiamy tych dwóch rodzin prostych
bez wyboru odpowiedniej rzeczywistej struktury.
Odwzorowanie $\tau$ jest gładkie,
ale nie nazywamy go holomorficznym:
podstawa $S^4$ nie została wyposażona
w strukturę zespoloną, podczas gdy
$\mathbb{CP}^3$ ma standardową taką strukturę.
\end{uwaga}

\section{Związki z algebrą liniową i dalsze kierunki}
\label{sec:lin32-wnioski}
Podstawowym obiektem obu konstrukcji jest
zespolona przestrzeń $\mathbb T\cong\C^4$
i jej rzutowanie $\mathbb{CP}^3$.
Relacja incydencji wybiera w niej
wykresy $W_X$ macierzy $2\times2$.
Struktura kwaternionowa wybiera zamiast
tego płaszczyzny niezmiennicze względem $J$.

Oba modele prowadzą do Grassmannianu
$\operatorname{Gr}_2(\C^4)$: jest on
przestrzenią \emph{wszystkich} prostych
rzutowych w $\mathbb{CP}^3$.
Jego obraz Plückera to kwadryka Kleina,
a rodzina $J$-niezmienniczych płaszczyzn
odpowiada rzeczywistej podrozmaitości
$\mathbb{HP}^1\cong S^4$.
Macierze rzędu jeden kodują z kolei
dwa typy maksymalnych kierunków
izotropowych, związanych ze spinorami.

Rozwinięcie tych idei do transformacji
Penrose'a, korespondencji Penrose'a--Warda
i analizy pól spinorowych wymaga teorii
snopów, kohomologii oraz koneksji.
Nie wynika ono automatycznie z dowiedzionych
tu elementarnych twierdzeń o incydencji
i fibracji.

\par\smallskip
{\footnotesize\noindent\emph{Dalsza lektura:}
M.~Atiyah, N.~Hitchin, I.~Singer,
\emph{Self-duality in four-dimensional Riemannian geometry};
przegląd teorii twistorów
\url{https://pmc.ncbi.nlm.nih.gov/articles/PMC5666237/};
opis kwaternionowy
\url{https://www.math.columbia.edu/~woit/twistorp1.pdf}.
W źródłach bywają używane odmienne konwencje
znaków i czynnika $i$ w równaniu incydencji.\par}
\clearpage
